% Part: sets-functions-relations
% Chapter: infinite
% Section: dedekind
%
\documentclass[../../../include/open-logic-section]{subfiles}

\begin{document}

\olfileid[pt]{sfr}{infinite}{dedekind}
\olsection{Álgebras de Dedekind}

Não queremos apenas que os números naturais sejam infinitos; queremos
que tenham certas propriedades (algébricas): precisam de comportar-se
bem sob adição, multiplicação, e assim por diante.

A ideia de Dedekind foi tomar como primitiva a noção de \emph{função
sucessor} e depois caracterizar os números como aqueles que têm as
seguintes propriedades:
\begin{enumerate}
	\item Existe um número, $0$, que não é sucessor de nenhum número
	\\O.s.\@, $0 \notin \ran{s}$
	\\O.s.\@, $\forall x\ s(x) \neq 0$
	\item Números distintos têm sucessores distintos
	\\O.s.\@, $s$ é !!a{injection}
	\\O.s.\@, $\forall x \forall y (s(x) = s(y) \lif x = y)$
	\item\ollabel{repeatedapplication} Todo o número se obtém de
	$0$ por aplicações repetidas da função sucessor.
\end{enumerate}
As duas primeiras condições são fáceis de lidar com, usando lógica de primeira
ordem (ver acima). Mas não podemos tratar \olref{repeatedapplication}
apenas com lógica de primeira ordem. O avanço de Dedekind foi reformular
a condição \olref{repeatedapplication}, em termos de teoria dos conjuntos,
assim:
\begin{enumerate}
	\item[3$'$.] Os números naturais são o menor conjunto que é
	\emph{fechado para a função sucessor}: isto é, se aplicarmos
	$s$ a qualquer !!{element} do conjunto, obtemos outro !!{element}
	do conjunto.
\end{enumerate}
Mas precisamos de explicitar isto com calma.

\begin{defn}\ollabel{Closure}
	Para qualquer função $f$, o conjunto $X$ é $f$-\emph{fechado} {sse}
	$(\forall x \in X)f(x) \in X$. Definimos agora, para qualquer $o$:
	$$\closureofunder{f}{o} = \bigcap\Setabs{X}{o \in X\text{ e }X
\text{ é $f$-fechado}}$$
\end{defn}

Logo $\closureofunder{f}{o}$ é a interseção de todos os conjuntos
$f$-fechados que têm $o$ como !!a{element}. Intuitivamente, então,
$\closureofunder{f}{o}$ é o conjunto $f$-fechado \emph{menor} que tem $o$
como !!a{element}. O resultado seguinte torna preciso esse pensamento intuitivo;
\begin{lem}\ollabel{closureproperties}
	Para qualquer função $f$ e qualquer $o \in A$:
	\begin{enumerate}
		\item\ollabel{closurehaselem} $o \in \closureofunder{f}{o}$; e
		\item\ollabel{closureclosed} $\closureofunder{f}{o}$ é $f$-fechado; e
		\item\ollabel{closuresmallest} se $X$ é $f$-fechado e $o \in
		X$, então $\closureofunder{f}{o} \subseteq X$
	\end{enumerate}
\end{lem}

\begin{proof}
Note-se que existe pelo menos um conjunto $f$-fechado com $o$ como !!a{element}, nomeadamente $\ran{f}\cup
\{o\}$. Logo $\closureofunder{f}{o}$, a interseção de \emph{todos} os
tais conjuntos, existe. Cumpre agora verificar
\olref{closurehaselem}--\olref{closuresmallest}.

Quanto a \olref{closurehaselem}: $o \in \closureofunder{f}{o}$ pois é uma
interseção de conjuntos que todos têm $o$ como !!a{element}.

Quanto a \olref{closureclosed}: suponhamos $x \in \closureofunder{f}{o}$. Se $o \in X$ e $X$ é $f$-fechado, então $x \in X$, e agora $f(x) \in X$ pois $X$ é
$f$-fechado. Logo $f(x) \in \closureofunder{f}{o}$.

Quanto a \olref{closuresmallest}: de modo geral, se $X
\in C$ então $\bigcap C \subseteq X$.
\end{proof}

Com isto, podemos dizer:

\begin{defn}
Uma \emph{álgebra de Dedekind} é um conjunto $A$ juntamente com uma função $f
\colon A \to A$ e algum $o \in A$ tal que:
	\begin{enumerate}
		\item \ollabel{ded:proper} $o \notin \ran{f}$
		\item \ollabel{ded:injection} $f$ é !!a{injection}
		\item \ollabel{ded:closure} $A = \closureofunder{f}{o}$
	\end{enumerate}
\end{defn}

Como $A = \closureofunder{f}{o}$, o resultado anterior nos diz que
$A$ é o menor conjunto $f$-fechado com $o$ como !!a{element}. Claramente uma
álgebra de Dedekind é infinita de Dedekind; basta olhar para as cláusulas
\olref{ded:proper} e \olref{ded:injection} da definição. Mas o fato mais
relevante é que qualquer conjunto infinito de Dedekind pode ser transformado
numa álgebra de Dedekind.

\begin{thm}\ollabel{thm:DedekindInfiniteAlgebra}
Se existe um conjunto infinito de Dedekind, então existe uma álgebra de Dedekind.
\end{thm}

\begin{proof}
Seja $D$ infinito de Dedekind. Logo existe uma injeção $g \colon D \to
D$ e um elemento $o  \in D \setminus \ran{g}$. Seja agora $A =
\closureofunder{g}{o}$; por \olref{closureproperties}, $A$ existe e $o \in A$. Seja $f =
\funrestrictionto{g}{A}$. Mostraremos que $A, f, o$ formam uma álgebra de
Dedekind.

Quanto a \olref{ded:proper}: $o \notin \ran{g}$ e $\ran{f}
\subseteq \ran{g}$, logo $o\notin \ran{f}$.

Quanto a \olref{ded:injection}: $g$ é injetiva em $D$; logo $f
\subseteq g$ tem de ser injetiva.

Quanto a \olref{ded:closure}: por \olref{closureproperties}, $A$ é $g$-fechado; \emph{a fortiori}, $A$ é $f$-fechado. Logo $\closureofunder{f}{o} \subseteq A$ por \olref{closureproperties}. Como também $\closureofunder{f}{o}$ é $f$-fechado e $f = \funrestrictionto{g}{A}$, segue-se que $\closureofunder{f}{o}$ é $g$-fechado. Logo $A \subseteq \closureofunder{f}{o}$ por \olref{closureproperties}.
\end{proof}

\end{document}
